\documentclass[11pt,a4paper]{article}
\usepackage[T1]{fontenc}
\usepackage[utf8]{inputenc}
\usepackage{amsmath,amssymb,amsthm}
\usepackage[margin=1in]{geometry}
\usepackage[colorlinks=true,linkcolor=blue,urlcolor=blue]{hyperref}
\theoremstyle{plain}
\newtheorem{theorem}{Theorem}
\newtheorem{lemma}[theorem]{Lemma}
\newtheorem{proposition}[theorem]{Proposition}
\newtheorem{corollary}[theorem]{Corollary}
\theoremstyle{definition}
\newtheorem{remark}[theorem]{Remark}

\title{Recovering the unwritten recurrences of the table OEIS A237340}
\author{Adrian Perez Fontelles\\ \small Independent researcher}
\date{2 September 2026}
\begin{document}
\maketitle

\begin{abstract}
OEIS A237340 is a two-dimensional table $T(n,k)$ whose empirical recurrences are, for several of
its lines, given only as an ORDER --- ``k=2: [order 25]'' and the like --- with the
coefficients in a linked file rather than in the entry. Those conjectures can still be settled,
and the recurrences recovered. A line of the table is a fixed-width or fixed-height array
count, hence a walk count on finitely many vertices, hence satisfies some monic recurrence of
order at most that number; Berlekamp--Massey on twice as many exact terms returns the MINIMAL
one. Where its order equals the order the entry states --- and where the same holds after the
first few terms are dropped --- any recurrence of that order the line satisfies has a
characteristic polynomial that is a multiple of the minimal one and of equal degree, hence is
that polynomial. This note settles column 2, column 3.
\end{abstract}

\noindent\small 2020 Mathematics Subject Classification. 05A15, 11B37, 15A18.\normalsize

\section{The table and the unwritten conjectures}

OEIS A237340 is ``T(n,k)=Number of (n+1)X(k+1) 0..3 arrays with the maximum minus the minimum of every 2X2 subblock equal.''. It is read by antidiagonals and begins
\[
256, 1980, 1980, 15082, 41220, 15082, 120808, 821234,\ \dots
\]
The entry carries, among its empirical claims:
\begin{quote}\small
k=2: [order 25]\\
k=3: [order 59]
\end{quote}
As of the ``Last modified'' line on the live entry (Jul 23 2025 10:06:42, revision 6) these are still
recorded as empirical, and the recurrences themselves are not in the entry text.

\section{Why a line of the table is a walk count}

Fix the index that the line fixes. The entry defines $T(n,k)$ as a number of arrays of a shape
determined by $n$ and $k$, subject to a condition local to a bounded window of consecutive
lines. Substituting the fixed index into the entry's own wording turns the definition into an
ordinary one-parameter array count, and for such a count the admissible windows are the
vertices of a finite digraph and an array is a walk. So the line is a walk count on $S$
vertices, and by the Cayley--Hamilton theorem it satisfies a monic integer recurrence of order
at most $S$.

\section{Recovering the recurrences}

\begin{lemma}\label{lem:bm}
A sequence known to satisfy some linear recurrence of order at most $S$ is determined, with its
minimal recurrence, by its first $2S$ terms; Berlekamp--Massey applied to them returns that
minimal recurrence exactly.
\end{lemma}

\subsection*{Column $k=2$}
Substituting the fixed index into the entry's wording gives an ordinary one-parameter array
count, a walk on $S=144$ vertices. Berlekamp--Massey on $2S=288$ exact terms
returns a minimal recurrence of order $25$ --- the order the entry states --- with
integer coefficients
\[
(c_1,\dots,c_{25})=(47,\ -506,\ -5350,\ 119418,\ -292491,\ -5769290,\ 38319506,\ 16591825,\ -782193687,\ 1731908806,\ 3853403204,\ -18163240306,\ 6781065371,\ 52909631318,\ -59067161032,\ -56301011344,\ 103756990512,\ 14034541408,\ -73613340160,\ 8403519488,\ 22531661824,\ -3926882304,\ -2687410176,\ 403046400,\ 86507520).
\]
so that $a(m)=\sum_{i=1}^{25}c_i\,a(m-i)$ for every $m$. Removing the first
$25+4$ terms and repeating gives order $25$ again, which is what the
identification below needs.

\subsection*{Column $k=3$}
Substituting the fixed index into the entry's wording gives an ordinary one-parameter array
count, a walk on $S=506$ vertices. Berlekamp--Massey on $2S=1012$ exact terms
returns a minimal recurrence of order $59$ --- the order the entry states --- with
integer coefficients
\[
(c_1,\dots,c_{59})=(96,\ -768,\ -144583,\ 2829202,\ 76621036,\ -2048554707,\ -17019483384,\ 683567240153,\ 1086744525507,\ -127743080457938,\ 217443164912333,\ 14645443896234312,\ -53088580372811046,\ -1088187773179562294,\ 5355964651067399703,\ 54511939490975873301,\ -323994533015450689437,\ -1896945537002544508359,\ 13079140648558270326783,\ 46839605666590539966105,\ -371664639002531175405612,\ -830972909954014455908246,\ 7681265156740366030560244,\ 10619637168892207960568400,\ -118011843089118668653251200,\ -96626201581965824262347184,\ 1368693452985749224973135952,\ 600052081413121464819195152,\ -12116333080185364478120896704,\ -2186663074827188730540497504,\ 82516013154375959175260463552,\ 587961512709604226591444608,\ -434597259495358647173627778304,\ 46416367635108081332516404480,\ 1775212858128677213203483541504,\ -318006915646874448230062952448,\ -5626288700514175098561773299712,\ 1255732865667653650897975547904,\ 13807215782241199385055922049024,\ -3442234102847103154187638046720,\ -26109888585139089543294480564224,\ 6937000551716746743500393627648,\ 37739819299829870946346945806336,\ -10522321819436677877472574570496,\ -41184702458694573813266893766656,\ 12098620590204231859319976755200,\ 33315037147413997943851162009600,\ -10479526682616061801699485941760,\ -19426519737898909383123347177472,\ 6698158915266764439922113773568,\ 7806837442494074020061471834112,\ -3038899735152044101439261245440,\ -1993371259803009208486629212160,\ 915937902523721933829420613632,\ 268029265847128836973032112128,\ -162462738472851709487669575680,\ -6837350022791405570680160256,\ 12671602253896835291550842880,\ -1605487506676779824877404160).
\]
so that $a(m)=\sum_{i=1}^{59}c_i\,a(m-i)$ for every $m$. Removing the first
$59+4$ terms and repeating gives order $59$ again, which is what the
identification below needs.

\begin{theorem}
For each line above, the displayed recurrence holds for every index, and it is the recurrence
the entry records.
\end{theorem}

\begin{proof}
It holds by Lemma~\ref{lem:bm}. For the identification, the entry asserts that the line
satisfies SOME recurrence of the stated order, possibly only from some index on. Let $q$ be its
characteristic polynomial and $q_{\min}$ the minimal polynomial of the tail. Then
$q_{\min}\mid q$, and the second Berlekamp--Massey run --- on the line with its first few
terms removed --- shows $\deg q_{\min}$ equals the stated order, which is $\deg q$. Both are
monic, so $q=q_{\min}$.
\end{proof}

\section{Verification}

Three checks.

First, each line's model was compared against the entry's own published values for that line,
taken from the antidiagonal DATA read in either orientation and from the ``Table starts''
block, and agrees at every available term. This is the only point at which reading the entry's
English enters, and a misreading gives different counts.

Second, each recovered recurrence was evaluated directly on the model's values, with no
Berlekamp--Massey involved, and holds at every index tested.

Third, each order was computed twice by independent routes --- modulo a $61$-bit prime, where
every intermediate is a machine word, and again in exact rational arithmetic --- and once more
on the line with its first few terms dropped, which is what the identification requires. All
agree.

\begin{thebibliography}{9}
\bibitem{oeis} The OEIS Foundation, \emph{The On-Line Encyclopedia of Integer
Sequences}, \url{https://oeis.org}, 2026. Sequence A237340.
\bibitem{massey} J.~L.~Massey, Shift-register synthesis and BCH decoding, \emph{IEEE
Trans. Inform. Theory} \textbf{15} (1969), 122--127.
\bibitem{stanley} R.~P.~Stanley, \emph{Enumerative Combinatorics}, Volume 1, 2nd ed.,
Cambridge University Press, 2011. (Transfer-matrix method, Section 4.7.)
\bibitem{horn} R.~A.~Horn and C.~R.~Johnson, \emph{Matrix Analysis}, 2nd ed., Cambridge
University Press, 2013.
\end{thebibliography}

\end{document}
